\documentclass[11pt,a4paper]{article}
\usepackage[T1]{fontenc}
\usepackage[utf8]{inputenc}
\usepackage[english]{babel}
\usepackage{amsmath,amssymb}
\usepackage[margin=2cm]{geometry}
\setlength{\parindent}{0pt}
\setlength{\emergencystretch}{3em}
\allowdisplaybreaks

\title{About \texttt{matrices\_beyond.tex}}
\author{}
\date{}
\begin{document}
\maketitle
This note explains what the test data in \texttt{matrices\_beyond.tex} is used for. That file is input for the engine in \texttt{main.c}; it is read by the test \texttt{test\_linear\_algebra.c} in this folder.

\begin{enumerate}
\item \textbf{Purpose.} The file holds 18 true statements of linear algebra that the domain does not read or decide yet. The test checks the other direction of soundness: the engine may say that a statement is not supported, but it must never declare a true statement false, and it must not claim to have decided what it cannot check.

\item \textbf{What is in it.} Entries with $\pi$ and with a square root (Vietri, exercises 1 and 7; even $\sqrt{4}$, which is a rational number written with a root); entries with a letter in a statement that does not place the letter in a set (Abbena, Fino and Gianella, chapter 5, exercises [11] and [12]; the determinant of a $2 \times 2$ matrix; the theorem of Binet for $2 \times 2$ matrices; a rotation matrix); a $1 \times 1$ matrix compared with a number; a determinant written with bars around a matrix; a power with a letter as exponent; the rank--nullity theorem on one matrix, with a kernel. And, with a letter placed in $\mathbb{R}$, what paragraph 32g does not treat: the rank of a matrix with a letter, an equality between matrices with a letter, the determinant of an inverse with a letter (Abbena, Fino and Gianella, exercise [3]), an entry divided by the letter, a Greek letter.

\item \textbf{Where they go.} The rank and the inverse with a letter are stage 3 of \texttt{handoff-algebra-lineare.md} (the cases of rule 18t); equalities between matrices with letters and letters without a set are stage 5 (identities between polynomials); roots are stage 6 (exact arithmetic with square roots in the checker); kernels and dimensions are stage 2. When a stage makes one of these statements provable, it moves to a file of statements that must be proved: the statement on the values of $k$ for which a $2 \times 2$ determinant vanishes, which was here, is now in \texttt{letter\_determinants.tex}.

\item \textbf{How to run it.} From the top folder of the repository:
\begin{verbatim}
cc -std=c11 -O2 -Wall -Wextra -Werror -o /tmp/test_linear_algebra \
    tests/linear_algebra/test_linear_algebra.c && /tmp/test_linear_algebra
\end{verbatim}

\item \textbf{Expected result.} The lines \texttt{PASS: none of the 18 true statements \dots{} is called false} and \texttt{PASS: 18 of 18 statements \dots{}, beyond the language, are left undecided}.
\end{enumerate}
\end{document}
